\documentclass[11pt, letterpaper]{article}
\usepackage[margin=0.45in]{geometry}
\usepackage[T1]{fontenc}
\usepackage{lmodern}
\usepackage{parskip}
\usepackage{hyperref}
\usepackage{microtype}
\usepackage{comment}
\usepackage{ifthen}
\usepackage{amssymb}
\usepackage{xcolor}

\newboolean{showcomments}
\setboolean{showcomments}{true}

% use \setboolean{showcomments}{false} anywhere in your document after
% importing this to disable the comments

\makeatletter
\newcommand{\mynote}[3]{%
  \ifthenelse{\boolean{showcomments}}{%
   \fbox{\bfseries\sffamily\scriptsize#1}%
   {\small$\blacktriangleright$\textsf{\emph{\color{#3}{#2}}}$\blacktriangleleft$}}%
  {%
   % these two lines ensure that there is no blank space inserted
   \@bsphack
   \@esphack
  }%
}
\makeatother

% One command per author:
\definecolor{asparagus}{rgb}{0.53, 0.66, 0.42}
%\newcommand{\jc}[1]{\mynote{Jeronimo}{#1}{asparagus}}



\newcommand{\SMG}[1]{\mynote{SMG}{#1}{red}}    %Girvin
\newcommand{\shraddha}[1]{\mynote{SS}{#1}{blue}}     %Singh







\begin{document}

\begin{center}
{\large \textit{Non-Abelian Composite Pulses in Quantum Phase Space}}\\[2pt]
Singh, Royer, Girvin
\end{center}

\vspace{4pt}

The theoretical language of non-Abelian quantum signal processing (NA-QSP) --- extending QSP to non-commuting control variables --- has been discussed in prior work, but no explicit non-Abelian composite pulse had ever been constructed. This paper provides the first: the Gaussian Controlled Rotation (GCR). The key physical insight is that non-commutativity, rather than being an obstacle to overcome, is a resource to exploit --- and exploiting it yields strictly better control sequences than standard Abelian approaches can provide. We then show that this single pulse sequence is sufficient to analytically match state-of-the-art numerical protocols for bosonic state preparation across all relevant state families, and to complete the analytical framework for universal control of GKP bosonic error-corrected qubits --- two domains where only numerical or partial analytical results previously existed. This central result and its applications are fully analytical with complete numerical verification.

Non-commuting control variables are a feature of quantum mechanics broadly, arising in superconducting circuits, trapped ions, and neutral atom arrays --- among the most actively developed experimental platforms today. The ability to synthesize nonlinear polynomials in non-commuting position and momentum operators directly enables Hamiltonian simulation of anharmonic oscillators, interacting bosonic models, spin-boson systems, and molecular dynamics on non-parabolic potential energy surfaces --- problems of direct interest across quantum chemistry, condensed matter, and high energy physics. Just as QSP was generalized to quantum singular value transformation (QSVT), unifying essentially all known qubit-only quantum algorithms, NA-QSP points toward a non-Abelian QSVT that may similarly unify hybrid CV-DV quantum algorithms --- an important open problem this paper explicitly motivates.

Section~II organizes CV-DV QSP protocols into a hierarchy by commutativity of control variables, situating GCR within the broader landscape. Section~III introduces GCR with an example targeted at coherent states, proves its correctness analytically, derives full error bounds, and benchmarks it against BB1, further introducing BB1(GCR) --- a concatenated construction whose non-unitary extension provides a natural seed for numerical optimization of the best available GKP readout circuit. This is the \textbf{first main result}: the first explicit non-Abelian composite pulse, achieving with two gates and 0.4\,$\mu$s circuit duration the same error scaling as the well-known Abelian BB1 sequence --- which requires four gates and 3.6\,$\mu$s --- while delivering a minimum of $10\times$ lower failure rate across all operating regimes.

Section~IV extends GCR to non-Gaussian states relevant for quantum error correction and quantum advantage with bosonic systems, providing the first fully analytical state preparation circuits for squeezed, cat, GKP, and Fock states using the phase-space instruction set --- the same used in a recent \textit{Nature Physics} experiment demonstrating fast universal control of a superconducting oscillator via numerically optimized circuits. That GCR performs on par with those numerical results while additionally enabling mid-circuit ancilla error detection --- a structural advantage no purely numerical scheme can guarantee --- is the \textbf{second main result}: a single composite pulse sequence suffices to match state-of-the-art machine learning based bosonic control across all relevant state families.

Section~V then uses GCR to yield the \textbf{third main result}: the complete analytical picture for universal control of GKP codes in superconducting circuits, directly building on the record-breaking beyond-break-even GKP experiments --- which demonstrated quantum error correction but lacked a complete framework for universal logical control --- without requiring any additional hardware resources. This includes \emph{piecewise gate teleportation}---to our knowledge the first high-fidelity, error-corrected single- and two-qubit logical gates for finite-energy GKP, surpassing all prior schemes even in the noiseless setting and autonomously protected against biased-noise ancilla errors as these experiments scale toward fault-tolerant operation. Section~VI explores the reach and limits of the phase-space instruction set toward universal state preparation, identifying concrete open problems for future work.

A general constructive theory of NA-QSP --- characterizing all polynomial transformations achievable with non-commuting control variables --- does not yet exist. GCR is the first explicit instance of what we expect to become a broad constructive theory of NA-QSP; by establishing that such non-Abelian composite pulses exist and are practically superior to Abelian ones, this work opens that program and delimits its open problems. Taken together, these results establish non-Abelian composite pulses as a practical and immediately deployable tool, with analytical structure, experimental readiness, and a clear theoretical horizon toward a full theory of NA-QSP and its extension to NA-QSVT --- at the intersection of quantum control, quantum error correction, and algorithm development.

\end{document}